\documentclass[11pt]{article}

% -------------------------------------------------
% Packages
% -------------------------------------------------
\usepackage{amsmath, amssymb}
\usepackage{graphicx}
\usepackage{booktabs}
\usepackage{siunitx}
\usepackage{caption}
\usepackage{subcaption}
\usepackage{float}
\usepackage{geometry}
\usepackage{enumitem}
\usepackage{hyperref}

\geometry{letterpaper, margin=1in}

\title{ESE 5720 Final Project}
\author{Krishna Chemudupati}

% -------------------------------------------------
\begin{document}
% -------------------------------------------------

\maketitle
\tableofcontents
\newpage

% -------------------------------------------------
% INTRODUCTION & LITERATURE REVIEW
% -------------------------------------------------
\section{Introduction and Literature Search}

\subsection{Motivation and Problem Statement}

High-speed optical links rely on a front-end that can convert very small, broadband photodiode currents into voltage signals large enough to drive subsequent baseband circuitry while preserving signal integrity. The key element in this front end is the transimpedance amplifier (TIA), which converts the photodiode current into a voltage with prescribed gain, bandwidth, and noise performance~\cite{course_handout}.

In this project, the objective is to design a wideband CMOS TIA in a 45\,nm process operating from a 1\,V supply. The TIA interfaces to a photodiode modeled (in small-signal) as a current source in parallel with parasitic capacitance $C_P \approx 100$\,fF, and it must drive a $600\,\Omega$ differential resistive load ($300\,\Omega$ from each output to ground). The required differential transimpedance gain is at least $150$\,k$\Omega$, with a minimum differential output swing of $0.4$\,V\textsubscript{pp}, an upper --3\,dB bandwidth exceeding 300\,MHz, and integrated input-referred current noise below 125\,nA\textsubscript{rms}, all under a total power budget below 10\,mW~\cite{course_handout}.

The parallel combination of the photodiode capacitance and the TIA’s input impedance sets a dominant pole at the front end. A high transimpedance implemented purely with a large resistor at the photodiode node would therefore severely limit bandwidth and increase thermal noise. Modern CMOS TIAs are designed to circumvent this gain--bandwidth--noise tradeoff through appropriate circuit topologies, feedback techniques, and device-level optimizations.

\subsection{Overview of TIA Topologies}

Transimpedance amplifiers (TIAs) are a critical front-end element in high-speed optical receivers. Their role is to convert the photocurrent generated by a photodiode into a voltage signal while preserving the bandwidth and minimizing the noise added in the process. TIA topologies must meet strict trade-offs between transimpedance gain, bandwidth, input impedance, noise performance, and power consumption, especially under low-voltage operation in advanced CMOS technologies. Architectures typically fall into one of the following categories: resistive feedback amplifiers, common-gate (CG) or regulated cascode (RGC) input stages, inverter-based topologies with active feedback, and distributed or inductively-peaked multistage designs.

\paragraph{Resistive Termination and Shunt-Feedback TIAs}

The most straightforward transimpedance amplifier topology is to terminate the photodiode with a passive resistor $R$ and sense the resulting voltage swing. The gain is given by Ohm's Law, but this approach creates a bandwidth-limiting pole at $f_{-3\text{dB}} = 1/(2\pi R(C_P + C_{\text{in}}))$, where $C_P$ is the parasitic capacitance of the photodiode and $C_{\text{in}}$ is the input capacitance of the sensing amplifier~\cite{compare}. As $R$ increases to improve gain, bandwidth rapidly decreases and thermal noise ($\sqrt{4kTR}$) increases, making this method unsuitable for high-speed designs.

To improve upon this, a high-gain amplifier with a feedback resistor $R_F$ can be used to implement a shunt-feedback TIA. This configuration virtually grounds the photodiode node, significantly reducing the input impedance and pushing the input pole outward. The closed-loop transimpedance is roughly equal to $R_F$, and the bandwidth is improved due to the negative feedback reducing the effect of $C_P$. However, real amplifiers have finite gain-bandwidth products and introduce their own noise and input capacitance, which limit overall performance. Amplifier noise becomes significant at low supply voltages, and loop stability must be verified across frequency and process corners~\cite{compare}.

Despite these challenges, several successful implementations have demonstrated high-speed performance. For instance, a 40\,Gb/s shunt-feedback TIA in 0.18\,$\mu$m CMOS achieved a 30.5\,GHz bandwidth and 51\,dB$\Omega$ gain by using cascaded common-source stages with $\pi$-type inductor peaking~\cite{tia40gbps}. This type of compensation helps extend bandwidth while maintaining flat gain across frequency.

\paragraph{Common-Gate and Regulated-Cascode Input Stages}

The common-gate (CG) configuration offers a different method of achieving low input impedance, where the source of a MOSFET receives the photodiode current and the gate is AC grounded. The input impedance is approximately $1/g_m$, and this directly lowers the RC time constant formed with $C_P$, pushing the input pole to higher frequency. This enables large bandwidths with relatively modest power consumption. However, CG TIAs suffer from high input-referred current noise since the full drain current noise of the transistor is transferred to the input~\cite{cmos_tia_ijvlsi}.

To mitigate these issues, regulated cascode (RGC) architectures insert a gain stage in a local feedback loop between the drain and gate of the input transistor. This increases the effective transconductance and decreases input impedance further. The RGC architecture has shown strong bandwidth extension and noise suppression capabilities. For example, an RGC TIA in 65\,nm CMOS achieved a bandwidth of 11.4\,GHz and a gain of 46\,dB$\Omega$ with just 23.9\,mW power dissipation, outperforming traditional shunt-feedback designs in both speed and energy per bit~\cite{rgc_mdpi}. The RGC structure is particularly attractive in low-voltage CMOS processes because it avoids stacking many transistors while maintaining performance.

\paragraph{Active Feedback and Inverter-Based TIAs}

In deep-submicron CMOS, supply voltages below 1\,V impose strict limitations on analog circuit headroom. Inverter-based TIAs have emerged as a promising alternative due to their simplicity and high transconductance when biased in the linear region. These designs typically include active (sometimes positive) feedback to lower input impedance and extend bandwidth~\cite{inverter_mdpi}.

However, active feedback loops introduce the risk of instability, especially when implemented with high-gain stages. Special care must be taken to control the gain-bandwidth product and ensure proper phase margin. For instance, loop compensation techniques or digitally tunable feedback paths may be used to guarantee stable operation across temperature and process variations~\cite{inverter_mdpi2}.

A representative implementation includes a 180\,nm CMOS inverter TIA that handles 1\,mA peak-to-peak input current by dynamically activating the feedback path only under high signal conditions. This configuration allows for automatic range extension without a dedicated limiting amplifier~\cite{inverter_mdpi3}. Another advanced design in 45\,nm CMOS targets LiDAR systems and uses dual current-mode feedback paths. It offers a tunable gain-bandwidth profile: 91\,dB$\Omega$ gain at 114\,MHz or 82\,dB$\Omega$ gain at 297\,MHz, while consuming only 245\,$\mu$W~\cite{inverter_45nm_bohrium}. This balance of low power, high transimpedance, and wide bandwidth makes inverter-based TIAs ideal for short-range optical sensing and battery-powered links.

\paragraph{Distributed, Inductively Peaked, and System-Level Designs}

At extremely high data rates (40\,Gb/s and beyond), the limitations of lumped-element TIAs motivate distributed amplifier architectures. These use a traveling-wave topology where gain stages are distributed along a terminated transmission line. Bandwidths exceeding 36\,GHz have been demonstrated in hybrid integrated receivers using CMOS TIAs co-designed with photodiodes and SiGe front-ends~\cite{receiver100gbps_optica}.

Inductively peaked multistage amplifiers are a more compact alternative, where on-chip spiral inductors are used to extend bandwidth by introducing a conjugate zero that cancels the input pole. These techniques are critical for ensuring flat gain and low group delay in 100\,Gb/s systems, and the TIA must be co-designed with equalizers and limiting amplifiers in mind. Packaging and photodiode bonding techniques (e.g., flip-chip, microbumps) also affect input capacitance and high-frequency behavior~\cite{compare}.

\paragraph{Low-Noise Combined TIA/LA (CTLA) Structures}

For systems that must detect both weak and strong optical pulses (e.g., LiDAR), TIAs with large dynamic range are essential. Combined transimpedance and limiting amplifier (CTLA) designs achieve this by fusing linear conversion and output clamping into a single stage. One such design in 180\,nm CMOS uses both a passive feedback resistor and active inverter feedback, along with an NMOS clamp to limit output swing~\cite{ctla_mdpi}. It achieves 88.8\,dB$\Omega$ transimpedance, 629\,MHz bandwidth, and supports input currents up to 1\,mA, making it well-suited for time-of-flight ranging and optical pulse detection.

\subsection{Implications for this Design}

The architectural review highlights several key principles for selecting a suitable TIA topology for implementation in a 45\,nm, low-voltage CMOS process. First and foremost, the input stage plays a critical role in defining both the bandwidth and the noise performance of the entire receiver front-end. Architectures such as the common-gate (CG), regulated cascode (RGC), and inverter-based feedback TIAs are particularly effective in reducing the input impedance presented to the photodiode. This impedance reduction suppresses the bandwidth-limiting pole formed with the photodiode's parasitic capacitance $C_P$, thus enabling high-speed operation without sacrificing gain. However, these benefits come with increased sensitivity to device sizing and bias point, since the input transistor noise (especially in CG and inverter structures) becomes the dominant contributor to input-referred current noise. Consequently, careful noise optimization is required to avoid violating the integrated noise constraint.

Second, the 1\,V supply constraint in 45\,nm CMOS technology imposes a hard limit on voltage headroom, which restricts transistor stacking in analog signal paths. Multi-stage cascode amplifiers with global feedback loops are often incompatible with these supply constraints. Instead, architectures based on local feedback (as seen in RGC and active-feedback inverter designs) are favored for their simplicity, superior power efficiency, and layout compatibility in deep submicron nodes.

Third, the output interface presents its own challenges. The TIA must drive a differential 600\,$\Omega$ load (300\,$\Omega$ to ground per side), which represents a relatively low-impedance termination. This mandates that the output stages of the TIA provide both high current drive capability and linearity while maintaining signal swing within the supply limits. To achieve this, gain is often partitioned across multiple stages: the initial stage prioritizes low input impedance and noise performance, while downstream amplifiers are optimized for transimpedance gain and load drive.

Fourth, the inclusion of any feedback mechanism---whether passive (shunt resistor around an op-amp) or active (e.g., an inverter-based loop)---requires robust stability analysis. The interaction between photodiode capacitance, on-chip parasitics, and feedback loop phase margin must be carefully simulated across process-voltage-temperature (PVT) corners to ensure that the TIA does not oscillate or exhibit peaking that violates noise or gain flatness specifications.

Based on these considerations, this project adopts a hierarchical multistage TIA architecture. The first stage will be a low input impedance current amplifier, realized using a CG or RGC structure. This ensures minimal bandwidth degradation due to the photodiode capacitance. The subsequent stages will be responsible for providing the majority of the transimpedance gain and for buffering the output to drive the 300\,$\Omega$ terminations. Where required, source-follower or resistively loaded differential amplifier stages will be introduced to balance bandwidth, gain, and output swing constraints.

The upcoming process characterization phase will provide the necessary device metrics---such as $g_m$, $r_o$, $C_{\text{gs}}$, $C_{\text{gd}}$, and unity-gain frequency $f_T$---under various width and bias current conditions. These parameters will be used in conjunction with small-signal models and noise analysis to derive sizing equations and initial biasing conditions. This foundation will support a structured design flow involving hand calculations, followed by Cadence simulations to validate AC gain, bandwidth, transient behavior, and input-referred noise against the project specifications.

\bibliographystyle{ieeetr}
\bibliography{tia_refs}


\newpage

% -------------------------------------------------
% PROCESS CHARACTERIZATION (PART A)
% -------------------------------------------------
\section{Process Characterization}

The design of any analog front-end begins with accurate process characterization of the target CMOS technology. For this project, we adopt the gpdk045 45\,nm CMOS process with a nominal supply voltage of 1\,V. The primary objective of this phase is to empirically extract the DC, small-signal, and high-frequency parameters of NMOS and PMOS devices across a range of channel widths. The simulation tool used is Cadence Virtuoso with Spectre.

\subsection{Simulation Procedure}

For both NMOS and PMOS devices, the following simulations are performed:

\begin{enumerate}
    \item \textbf{$I_D$ vs.\ $V_{DS}$ Characteristics:} Sweep $V_{DS}$ from 0 to 1\,V (10 uniform points) for a fixed $V_{GS}$ in saturation. Device length is fixed at $L = 45$\,nm, and widths of $W = \{0.5\,\mu\mathrm{m}, 1\,\mu\mathrm{m}, 2\,\mu\mathrm{m}, 4\,\mu\mathrm{m}, 8\,\mu\mathrm{m} \}$ are used to observe output resistance and saturation behavior.
    
    \item \textbf{$V_{th}$ vs.\ $I_D$:} Using DC sweeps of $V_{GS}$, extract threshold voltage as a function of drain current for the same widths. Threshold is determined using the constant-current method (e.g., $I_D = 100\,\mathrm{nA} \cdot W/L$), providing insight into subthreshold and moderate inversion regimes.
    
    \item \textbf{$g_m$ vs.\ $I_D$:} Transconductance is extracted from small-signal AC simulations or derived from the slope of the $I_D$–$V_{GS}$ curve. This plot is essential to determine $g_m/I_D$ efficiency across inversion levels. Device sizing for optimal power efficiency will reference this plot.
    
    \item \textbf{$C_{\text{gs}}$ and $C_{\text{gd}}$ vs.\ $W$:} Gate capacitances are extracted at 1\,V bias and plotted versus width from 0.4\,$\mu$m to 10\,$\mu$m. This data informs input capacitance estimation and photodiode pole placement.
    
    \item \textbf{Unity Gain Frequency ($f_T$):} Extracted from AC simulation using the relation $f_T = g_m/(2\pi(C_{\text{gs}} + C_{\text{gd}}))$. Sweeps are performed across multiple bias currents to observe how $f_T$ varies with $I_D$. This metric determines the maximum achievable bandwidth and limits the topological complexity of high-speed amplifiers.
\end{enumerate}

\subsection{Data Usage and Validation}

All curves will be presented for both NMOS and PMOS transistors, enabling design symmetry in differential stages. The extracted data is used to:

\begin{itemize}
    \item Validate hand calculations of gain, input impedance, and bandwidth.
    \item Estimate noise and parasitic poles.
    \item Calibrate device models for schematic-level simulations.
\end{itemize}

Where discrepancies are observed between simulated data and textbook equations (e.g., square-law or EKV model predictions), these will be explained in the report to account for velocity saturation, channel length modulation, and short-channel effects.

The outcome of this process is a set of device curves and operating point tables that support all subsequent design stages, ensuring that transistor sizing aligns with bandwidth, gain, noise, and power specifications.


\newpage

\subsection{$I_D$ vs $V_{DS}$}

To begin the process characterization, the drain--current versus drain--source
voltage characteristics of minimum-length MOSFETs were measured for a range of
device widths. All NMOS transistors were implemented using the $45\,\mathrm{nm}$
process with channel length $L = 45\,\mathrm{nm}$ and widths
$W \in \{0.5\,\mu\mathrm{m},\, 1\,\mu\mathrm{m},\, 2\,\mu\mathrm{m},\, 4\,\mu\mathrm{m},\, 8\,\mu\mathrm{m}\}$.
For each device, the gate--source voltage was held at $V_{GS} = 1\,\mathrm{V}$,
and the drain--source voltage $V_{DS}$ was swept from $0$ to $1\,\mathrm{V}$ in uniform
steps. The drain current $I_D$ at the drain terminal was recorded as the primary
output of the DC analysis. An identical procedure was used for PMOS devices,
with the source and bulk tied to $V_{DD} = 1\,\mathrm{V}$ and the gate held at
ground, resulting in a constant $V_{SG} = 1\,\mathrm{V}$. In both cases, sweeping
the device width produces a corresponding family of $I_D$--$V_{DS}$ curves.

The resulting characteristics reflect the expected long-channel trends discussed
in the course notes, even though the $45\,\mathrm{nm}$ devices depart from ideal
behavior. At low drain--source voltages ($V_{DS} < V_{GS} - V_{TH}$), the NMOS
operates in the triode region and the drain current increases approximately
linearly with $V_{DS}$, with slope proportional to $W/L$. As $V_{DS}$ approaches
the saturation boundary, $V_{DS,\mathrm{sat}} \approx V_{GS} - |V_{TH}|$, the
curves transition from a steep linear rise to a more gradual increase, consistent
with the onset of velocity saturation and channel-length modulation. Using the
operating-point data from the simulations, we extract threshold voltages of
approximately $V_{TH,n} \approx 610\,\mathrm{mV}$ for the NMOS devices and
$|V_{TH,p}| \approx 530\,\mathrm{mV}$ for the PMOS devices. These values place
the expected saturation voltages near
\[
  V_{DS,\mathrm{sat}} \approx V_{GS} - V_{TH,n} \approx 0.39\,\mathrm{V}
  \quad\text{and}\quad
  V_{SD,\mathrm{sat}} \approx V_{SG} - |V_{TH,p}| \approx 0.47\,\mathrm{V}.
\]

For $V_{DS}$ beyond the saturation boundary, the drain current continues to rise
slowly with $V_{DS}$ due to channel-length modulation, again with magnitude
scaling roughly in proportion to the device width. Across all widths, the
measured characteristics exhibit the expected behavior: nearly zero current at 
low $V_{DS}$, a rapid increase through the triode region, and a flattened slope
once the devices enter saturation. These observations confirm consistency with
standard MOSFET operation and form the basis for the subsequent extraction of
threshold voltage, transconductance, and $g_m/I_D$ characteristics used later in
the TIA design analysis.

\newpage

\subsubsection{NMOS Test Setup and Results}

\begin{figure}[H]
    \centering
    \includegraphics[width=0.9\linewidth]{writeup//plots/nmos_id_vds_tb.png}
    \caption{NMOS testbench setup}
\end{figure}

\begin{figure}[H]
    \centering
    \includegraphics[width=0.9\linewidth]{writeup//plots/nmos_id_vds_plot_ss.png}
    \caption{Results showing $I_{D, max}$ for each width}
\end{figure}

\newpage

\subsubsection{PMOS Test Setup and Results}

\begin{figure}[H]
    \centering
    \includegraphics[width=0.9\linewidth]{writeup//plots/pmos_id_vds_tb.png}
    \caption{NMOS testbench setup}
\end{figure}

\begin{figure}[H]
    \centering
    \includegraphics[width=0.9\linewidth]{writeup//plots/pmos_id_vds_plot_ss.png}
    \caption{Results showing $I_{D, max}$ for each width}
\end{figure}

\newpage

\subsubsection{Summary Table}

\begin{table}[h!]
\centering


\begin{tabular}{c c c}
\hline
\textbf{Width $W$ ($\mu$m)} & \textbf{$I_{D,\max}$ NMOS (mA)} & \textbf{$I_{D,\max}$ PMOS (mA)} \\
\hline
0.5 & 0.1807 & $-0.12635$ \\
1.0 & 0.3503 & $-0.25455$ \\
2.0 & 0.6775 & $-0.50322$ \\
4.0 & 1.319  & $-1.00010$ \\
8.0 & 2.587  & $-1.99110$ \\
\hline
\end{tabular}
\vspace{6pt}
\caption{Summary of $I_{D,\max}$ for NMOS and PMOS at $V_{DS}=1\,\mathrm{V}$ and $V_{GS}=V_{SG}=1\,\mathrm{V}$}
\end{table}



\newpage

\subsection{$V_{TH}$ vs $I_D$}

In the second characterization step, the threshold voltage of minimum-length
devices was evaluated as a function of drain current for several widths. Unlike
the diode-connected configuration sometimes used for $V_{TH}$ extraction, the
project specifications require that this characterization be performed under the
same bias conditions as the $I_D$--$V_{DS}$ measurements from Section~2.1.
Accordingly, each NMOS transistor was biased with $V_{DS}=1\,\mathrm{V}$, while
the gate--source voltage was swept from $0$ to $1\,\mathrm{V}$ in a DC analysis.
This sweep drives the drain current from the subthreshold region up to the
$I_{D,\max}$ values obtained previously (ranging from about $0.18\,\mathrm{mA}$
for $W=0.5\,\mu\mathrm{m}$ up to $2.59\,\mathrm{mA}$ for $W=8\,\mu\mathrm{m}$).
For each bias point, the transistor’s built-in operating-point parameter
\texttt{vth} was recorded, and the resulting threshold voltage was plotted
against the corresponding drain current for each width.

An analogous procedure was used for the PMOS devices. With the source and bulk
tied to $V_{DD}=1\,\mathrm{V}$ and the gate held at ground, a sweep of the
gate--source voltage magnitude $V_{SG}$ from $0$ to $1\,\mathrm{V}$ generates
drain currents up to approximately $2\,\mathrm{mA}$ for the widest device.
Again, the model’s \texttt{vth} parameter was extracted at each operating point
and plotted versus the magnitude of $I_{SD}$.

From the long-channel strong-inversion relations,
\[
  I_D \approx \frac{1}{2}\mu_n C_{\mathrm{ox}} \frac{W}{L}(V_{GS}-V_{TH})^2
\]
with fixed $V_{DS} \ge V_{GS}-V_{TH}$, the gate overdrive $V_{GS}-V_{TH}$ is
expected to vary as $\sqrt{I_D}$, while $V_{TH}$ itself remains approximately
constant. In short-channel technologies such as $45\,\mathrm{nm}$, however,
effects including drain-induced barrier lowering (DIBL), mobility degradation,
and channel-length modulation introduce a weak dependence of $V_{TH}$ on both
drain current and device geometry.

The simulated $V_{TH}$--$I_D$ characteristics for the NMOS devices exhibit the
expected behavior. Across the full current range, the extracted threshold values
cluster tightly around $V_{TH,n} \approx 610\,\mathrm{mV}$, with only a few
millivolts of variation as $I_D$ increases from the microampere range to several
milliamperes. Dependence on device width is likewise minimal: wider devices show
nearly the same threshold as narrower ones, aside from a small systematic shift
at the highest currents due to short-channel effects. The PMOS devices show
similar behavior, with extracted thresholds remaining close to
$|V_{TH,p}| \approx 530\,\mathrm{mV}$ throughout the sweep and only mild
broadening at large $I_{SD}$.

These results confirm that, over the bias region relevant to the subsequent TIA
hand calculations, both NMOS and PMOS devices may be modeled as having nearly
constant thresholds in the $0.5$--$0.65\,\mathrm{V}$ range. Second-order
short-channel effects introduce only modest deviations, and the extracted
$V_{TH}$ values provide a consistent foundation for the later $g_m$,
$g_m/I_D$, and device sizing analyses.

\newpage

\subsubsection{NMOS Test Setup and Results}

\begin{figure}[H]
    \centering
    \includegraphics[width=1\linewidth]{writeup//plots/nmos_vth_tb.png}
    \caption{}
\end{figure}

\begin{figure}[H]
    \centering
    \includegraphics[width=1\linewidth]{writeup//plots/nmos_vth_id_plot_ss.png}
    \caption{}
\end{figure}

\newpage

\subsubsection{PMOS Test Setup and Results}

\begin{figure}[H]
    \centering
    \includegraphics[width=1\linewidth]{writeup//plots/pmos_vth_tb.png}
    \caption{}
\end{figure}

\begin{figure}[H]
    \centering
    \includegraphics[width=1\linewidth]{writeup//plots/pmos_vth_id_plot_ss.png}
    \caption{}
\end{figure}

\newpage

\subsection{$g_m$ vs $I_D$}

In the third characterization step, the small–signal transconductance $g_m$ of
minimum–length MOSFETs was evaluated as a function of drain current. For this
measurement, the same diode–connected testbench topology used in the previous
section was employed. Each NMOS device with $L = 45\,\mathrm{nm}$ and widths
$W = \{0.5,\,1,\,2,\,4,\,8\}\,\mu\mathrm{m}$ was connected with its gate and drain
shorted, source and bulk tied to ground, and a DC current source applied from
drain to supply. A DC analysis was run while sweeping the imposed drain current
$I_D$ from $0$ up to a value slightly larger than the $I_{D,\max}$ found in the
$I_D$–$V_{DS}$ characterization (approximately $4\,\mathrm{mA}$ in 100 points).
At each operating point, the simulator’s operating–point parameter \texttt{gm}
was saved, and the resulting $g_m$–$I_D$ characteristics were plotted for all
device widths.

An analogous setup was used for the PMOS devices. With the source and bulk tied
to $V_{DD}=1\,\mathrm{V}$, each PMOS transistor was diode–connected by shorting
its gate and drain, and a DC source current $I_{SD}$ was swept from source to
$V_{DD}$. The current range was chosen to extend slightly above the
corresponding $I_{SD,\max}$ values from the PMOS $I_{SD}$–$V_{SD}$ curves,
resulting in a maximum magnitude of approximately $3\,\mathrm{mA}$. The simulator
again reported $g_m$ at each bias point, enabling the $g_m$–$I_{SD}$
characteristics to be compared across widths. From the long–channel strong–inversion model,
\[
  I_D \approx \frac{1}{2}\mu C_{\mathrm{ox}}\frac{W}{L}(V_{GS}-V_{TH})^2(1+\lambda V_{DS}),
\]
and with $V_{DS}=V_{GS}$ in a diode–connected configuration, differentiation
with respect to $V_{GS}$ yields
\[
  g_m = \frac{\partial I_D}{\partial V_{GS}}
  \approx \mu C_{\mathrm{ox}}\frac{W}{L}(V_{GS}-V_{TH})
  \approx \frac{2I_D}{V_{GS}-V_{TH}}.
\]
Eliminating the gate overdrive gives the familiar square–law scaling
\[
  g_m \approx \sqrt{\frac{2\mu C_{\mathrm{ox}}}{L}}\,\sqrt{W I_D},
\]
indicating that $g_m$ should grow approximately as $\sqrt{I_D}$ for a given
width, and increase proportionally to $\sqrt{W}$ at any fixed current.

The simulated curves in Figures 10 and 12 follow these expectations. For each device width, $g_m$ rises monotonically with increasing drain current and shows a sublinear trend at higher currents, consistent with mobility degradation and velocity saturation in the $45\,\mathrm{nm}$ technology. Wider devices exhibit larger transconductance at a given current, and the vertical spacing between curves reflects the predicted $\sqrt{W}$ dependence. As $I_D$ increases into the milliampere range, short–channel effects cause $g_m$ to deviate slightly from the ideal $\sqrt{W I_D}$ behavior, leading to a gentle flattening of the curves.

These $g_m$–$I_D$ characteristics provide valuable design insight. Because
transconductance is a key parameter in determining gain and bandwidth of the TIA
in later sections, the extracted curves allow rapid selection of the necessary
bias current for each device width. The predictable scaling with $\sqrt{W}$ and
the smooth increase of $g_m$ with $I_D$ form a reliable foundation for the
device–level sizing and biasing decisions used in the subsequent design stages.

\newpage

\subsubsection{NMOS Test Setup and Results}

\begin{figure}[H]
    \centering
    \includegraphics[width=1\linewidth]{writeup//plots/nmos_gm_tb.png}
    \caption{}
\end{figure}

\begin{figure}[H]
    \centering
    \includegraphics[width=1\linewidth]{writeup//plots/nmos_gm_id_plot_ss.png}
    \caption{}
\end{figure}

\newpage

\subsubsection{PMOS Test Setup and Results}

\begin{figure}
    \centering
    \includegraphics[width=1\linewidth]{writeup//plots/pmos_gm_tb.png}
    \caption{}
\end{figure}

\begin{figure}[H]
    \centering
    \includegraphics[width=1\linewidth]{writeup//plots/pmos_gm_id_plot_ss.png}
    \caption{}
\end{figure}

\newpage

\subsection{$C_{gs}$ and $C_{gd}$ vs $I_D$}

The final device characterization step investigates how the intrinsic gate
capacitances $C_{gs}$ and $C_{gd}$ scale with transistor width. For this
measurement, each MOSFET was placed in a diode–connected configuration by
shorting the gate and drain terminals, while the source and body were tied to
ground. This bias forces the device into strong inversion for sufficiently large
widths and establishes a well–defined small–signal operating point from which
the intrinsic capacitances can be extracted. A DC operating–point analysis was
run for each width, and the simulator’s operating–point quantities
\texttt{cgs} and \texttt{cgd} were recorded. Width was swept across the range
$0.4\,\mu\mathrm{m} \le W \le 10\,\mu\mathrm{m}$ using ten evenly spaced points.

From first–order MOS capacitor theory, the gate–source capacitance in
strong inversion is dominated by the channel–charge sharing between source and
drain, giving approximately
\[
  C_{gs} \approx \frac{2}{3} W L C_{\mathrm{ox}},
\qquad
  C_{gd} \approx 0
\]
for long–channel devices biased with $V_{DS}=V_{GS}$. In short–channel devices,
fringing fields, overlap capacitances, and velocity–saturation effects introduce
additional contributions. As a result, $C_{gs}$ is expected to scale nearly
linearly with width, while $C_{gd}$ should remain roughly width–independent
aside from gate–drain overlap and short–channel contributions.

The simulated results in Figures 14 and 16 follow these expectations. For both NMOS and PMOS devices, $C_{gs}$ increases almost linearly with width over the entire range from $0.4$ to $10\,\mu\mathrm{m}$, reflecting the proportional increase in gate area. The gate–drain capacitance $C_{gd}$ remains significantly smaller than $C_{gs}$ and varies only weakly with width, with slight deviations at large $W$ due to increased overlap and short–channel effects. The NMOS and PMOS capacitances show similar trends in magnitude, with differences arising from the underlying process parameters and carrier mobilities.

These capacitance–width characteristics provide the values needed for estimating
input node bandwidth and for developing the small–signal models used in the TIA
design that follows. In particular, the near–linear growth of $C_{gs}$ with $W$
is a key consideration when selecting device widths to balance gain, bandwidth,
and noise in later stages of the design.

\newpage

\subsubsection{NMOS Test Setup and Results}

\begin{figure}[H]
    \centering
    \includegraphics[width=0.9\linewidth]{writeup//plots/nmos_c_tb.png}
    \caption{}
\end{figure}

\begin{figure}[H]
    \centering
    \includegraphics[width=0.9\linewidth]{writeup//plots/nmos_caps_plots_ss.png}
    \caption{}
\end{figure}

\newpage

\subsubsection{PMOS Test Setup and Results}

\begin{figure}[H]
    \centering
    \includegraphics[width=1\linewidth]{writeup//plots/pmos_c_tb.png}
    \caption{}
\end{figure}

\begin{figure}[H]
    \centering
    \includegraphics[width=1\linewidth]{writeup//plots/pmos_caps_plots_ss.png}
    \caption{}
\end{figure}

\newpage

\subsection{Process $f$T}

The unity–gain frequency $f_T$ of a MOSFET is defined as the frequency at which
the short–circuit current gain of the device falls to unity, that is, the
magnitude of the small–signal drain current equals the magnitude of the
small–signal gate current. A convenient way to estimate $f_T$ is to model the
transistor as a transconductance $g_m$ driving a lumped capacitance at the gate
equal to the total input capacitance
\[
  C_L \approx C_{gs} + C_{gd}.
\]
In this simple one–pole model the magnitude of the current gain is
\[
  \left|\frac{i_d}{i_g}\right| \approx \frac{g_m}{\omega C_L},
\]
so the unity–gain condition $\lvert i_d / i_g \rvert = 1$ occurs at
\begin{equation}
  f_T \approx \frac{g_m}{2\pi C_L}.
  \label{eq:fT_def}
\end{equation}

Using the data from the $g_m$–$I_D$ and capacitance characterizations, we can
estimate the order of magnitude of $f_T$ for this $45\,\mathrm{nm}$ process. For
a representative minimum–length NMOS device with width
$W = 2\,\mu\mathrm{m}$ biased at $I_D \approx 0.4\,\mathrm{mA}$, the
$g_m$–$I_D$ curves give a transconductance of approximately
\[
  g_m \approx 3\,\mathrm{mS}.
\]
From the $C_{gs}$ and $C_{gd}$ versus width results, the total gate
capacitance for a device in this size range is on the order of a few
femtofarads; a reasonable range for the combined value is
\[
  C_L = C_{gs} + C_{gd} \approx 1.6\text{–}1.8\,\mathrm{fF}.
\]
Substituting these representative values into \eqref{eq:fT_def} gives
\[
  f_T \sim
  \frac{3\times10^{-3}}
       {2\pi (1.6\text{–}1.8)\times10^{-15}}
  \approx (2.7\text{–}3.0)\times10^{11}\,\mathrm{Hz},
\]
so the intrinsic unity–gain frequency of a minimum–length NMOS device in this
bias range is roughly
\[
  f_T \approx 3\times10^{11}\,\mathrm{Hz}
  \quad (\text{about } 300\,\mathrm{GHz}).
\]
This is comfortably higher than the upper frequency of interest for the TIA
design (around $10\,\mathrm{GHz}$), indicating that transit–time limits of the
individual devices should not be the dominant bottleneck for bandwidth.

Equation \eqref{eq:fT_def} also makes the dependence of $f_T$ on drain current
clear. In strong inversion, the transconductance can be written approximately as
\[
  g_m \approx \frac{2I_D}{V_{OV}},
\]
where $V_{OV} = V_{GS} - V_{TH}$ is the overdrive voltage. Over the current
range used in the preceding characterizations, the total input capacitance
$C_L$ is dominated by geometric and overlap components and therefore varies only
weakly with bias, whereas $g_m$ increases significantly with $I_D$. As a
result, within this operating region $f_T$ is roughly proportional to $g_m$ and
hence increases with drain current. At very large current densities mobility
degradation and velocity saturation cause $g_m$ to grow more slowly with $I_D$,
so the rate of increase of $f_T$ eventually tapers off, but over the bias range
relevant to this project the approximation that $f_T$ increases with $I_D$ is
adequate.

\newpage

\end{document}
